\chapter{Pushing to the Browser}

The screen backend now exists and holds a session with the plant. The other half
of its job is the browser. An operator screen needs a fresh picture of the plant
about twice a second: 41 measurements, 12 valve positions and whether the plant
has shut down. The browser has to receive that picture without asking for it,
keep receiving it when the network is slow, and know when the plant behind it is
gone.

Chapter~2 left one question open: how a backend pushes to a browser. The answer
took the web fifteen years.

\section{Asking again}

The web was built on one exchange: the browser asks, the server answers, and
the conversation is over. The server has no way to speak first. So the simplest
live screen is the one from Chapter~1, moved to the browser: ask every half
second.

\begin{figure}[htbp]
\centering
\begin{tikzpicture}[x=1cm, y=-0.62cm]
  \node[hf/note, font=\sffamily\footnotesize\bfseries] at (2.5,-1.3) {polling};
  \node[hf/reader, minimum width=18mm] at (0,0) {browser};
  \node[hf/reader, minimum width=18mm] at (5,0) {server};
  \draw[hf/lifeline] (0,0.7) -- (0,7.4);
  \draw[hf/lifeline] (5,0.7) -- (5,7.4);
  \foreach \y in {1.4,3.4,5.4} {
    \draw[hf/ask] (0,\y) -- node[hf/note, above] {GET} (5,\y);
    \draw[hf/answer] (5,\y+0.9) -- node[hf/note, above] {headers + 53 numbers} (0,\y+0.9);
  }
  \begin{scope}[xshift=7.5cm]
    \node[hf/note, font=\sffamily\footnotesize\bfseries] at (2.5,-1.3) {push};
    \node[hf/reader, minimum width=18mm] at (0,0) {browser};
    \node[hf/reader, minimum width=18mm] at (5,0) {server};
    \draw[hf/lifeline] (0,0.7) -- (0,7.4);
    \draw[hf/lifeline] (5,0.7) -- (5,7.4);
    \draw[hf/ask] (0,1.4) -- node[hf/note, above] {open} (5,1.4);
    \foreach \y in {2.9,4.4,5.9} \draw[hf/push] (5,\y) -- node[hf/note, above] {53 numbers} (0,\y);
  \end{scope}
\end{tikzpicture}
\caption{The same three updates. Polling pays a full request and response for
each one, per browser. Push pays for one connection and then sends only the
data.}
\end{figure}

Polling works, and its cost is easy to state. Each open tab makes two requests
a second. Ten tabs make twenty. Every answer carries a full set of HTTP headers,
often larger than the 53 numbers inside it, and many of those answers say
nothing new. The server does work in proportion to the number of tabs, not to
the number of changes in the plant.

\begin{history}{1999--2006 \textperiodcentered{} Ajax and Comet}
In 1999, the team building Outlook Web Access at Microsoft needed a mail page
that could fetch new messages without reloading. They added an object to
Internet Explorer 5 that let a script make an HTTP request in the background.
Other browsers copied it as \texttt{XMLHttpRequest}. In February 2005, Jesse
James Garrett gave the style of application built on it a name, \textbf{Ajax},
and Google Maps and Gmail made it famous.

Ajax still asked first. In March 2006, Alex Russell named the set of tricks
people used to make the server speak first \textbf{Comet}. The most common one
was the \emph{long poll}: the browser asks, the server holds the request open
until it has something to say, answers, and the browser immediately asks again.
It worked everywhere and it was a hack: one request per message, and a server
holding thousands of requests that were really waiting connections.
\end{history}

\section{A connection that stays open}

\begin{history}{2004--2011 \textperiodcentered{} Server-Sent Events and WebSocket}
Ian Hickson was the editor of the document that describes how browsers must
behave, the one that later got the name HTML5. He worked on the problem from
two sides. The first answer, sketched in 2004 and shipped by the Opera browser
in 2006, was \textbf{Server-Sent Events}: the browser makes one ordinary HTTP
request, and the server never finishes the response, writing one event after
another into it as plain text. It only goes one way, server to browser.

The second was a real two-way connection, where either side can send at any
moment. The draft had a slot reserved for it under a provisional name,
\texttt{TCPConnection}, because the idea was to give a page something close to
a raw network connection, and TCP is the protocol that carries almost every
connection on the internet. In June 2008, a developer named Michael Carter
pushed the discussion forward on the group's mailing list and chat channel. The
design took shape in a few weeks, and he and Hickson named it
\textbf{WebSocket}. Chrome shipped it at the end of 2009. Because it was a
network protocol and not just a browser feature, it then went to the body that
standardizes internet protocols, where it was rewritten to close security
holes, and was published as RFC~6455 in December 2011.
\end{history}

\begin{question}{Who decides how the web works?}
Three groups, and the history above passes through all of them.

The \textbf{W3C}, the World Wide Web Consortium, was founded by Tim
Berners-Lee in 1994 to standardize the web he had invented. By the early 2000s
it was betting on a stricter successor to HTML that browsers weren't
interested in. In June 2004, people from Apple, Mozilla and Opera left to
write the next HTML themselves, as the \textbf{WHATWG} (Web Hypertext
Application Technology Working Group). Hickson was its editor, and its document
became HTML5. In 2019 the W3C agreed that the WHATWG's version is the official
one. HTML and the rules for what a browser does are theirs.

The \textbf{IETF}, the Internet Engineering Task Force, has met since 1986 and
standardizes the protocols underneath: TCP, HTTP, and anything else that two
programs say to each other over the network. It decides by ``rough consensus
and running code'', in David Clark's 1992 phrase. A finished IETF standard is
published as an \textbf{RFC}, a numbered document that never changes; a
correction gets a new number. The name means \emph{Request for Comments}, and
it's a leftover from April 1969, when Steve Crocker, a graduate student, wrote
RFC~1 about the ARPANET and chose a title humble enough not to offend anyone
senior. The name stuck to documents that became the law of the internet.

So WebSocket has two halves: the browser's side, in the WHATWG's HTML, and the
protocol on the wire, in an IETF RFC.
\end{question}

A WebSocket starts life as an HTTP request, so it passes through the same ports,
proxies and load balancers as everything else on the web. The request asks the
server to stop speaking HTTP on this connection:

\begin{figure}[htbp]
\centering
\begin{tikzpicture}[x=1cm, y=-0.7cm]
  \node[hf/reader, minimum width=20mm] at (0,0) {browser};
  \node[hf/reader, minimum width=20mm] at (9,0) {server};
  \draw[hf/lifeline] (0,0.7) -- (0,7.8);
  \draw[hf/lifeline] (9,0.7) -- (9,7.8);
  \draw[hf/ask] (0,1.6) -- node[hf/note, above, align=center]
    {\texttt{GET /stream}\enspace\texttt{Upgrade: websocket}\\\texttt{Sec-WebSocket-Key: dGhlIHNhbXBsZSBub25jZQ==}} (9,1.6);
  \draw[hf/answer] (9,3.4) -- node[hf/note, above, align=center]
    {\texttt{101 Switching Protocols}\\\texttt{Sec-WebSocket-Accept: s3pPLMBiTxaQ9kYGzzhZRbK+xOo=}} (0,3.4);
  \node[hf/note, anchor=west] at (9.2,4.4) {no more HTTP};
  \draw[black!25, dashed] (-0.5,4.4) -- (9,4.4);
  \foreach \y in {5.2,6.0,6.8} \draw[hf/push] (9,\y) -- (0,\y);
  \node[hf/note, anchor=east] at (-0.2,6.0) {frames};
  \draw[hf/ask] (0,7.4) -- node[hf/note, above] {a frame the other way, any time} (9,7.4);
\end{tikzpicture}
\caption{The WebSocket handshake. One HTTP exchange, and then the same TCP
connection carries frames in both directions.}
\end{figure}

Look at the two odd lines in the figure, the \texttt{Key} going out and the
\texttt{Accept} coming back. They are a challenge and its answer. Before the
browser stops speaking HTTP, it wants proof that the other side read this
request and really speaks WebSocket. An ordinary web server that has never
heard of WebSocket shouldn't be able to say yes by accident.

So the browser sends a challenge: a random value, new for every connection.
The server has to transform it with a rule that every WebSocket server knows:
glue a fixed string, written once into RFC~6455, to the end of the challenge,
hash the result, and send it back. The browser does the same computation on its
side. If the two results match, the connection goes on. If they don't, the
browser closes it. A server that doesn't know the rule can't produce the right
answer, and an answer saved from an earlier connection is useless, because the
challenge is different every time.

The whole rule is one line of code. Here it is, run on the sample challenge
printed in the RFC, the same one in the figure:

\begin{code}
import { createHash } from 'node:crypto';

const key = 'dGhlIHNhbXBsZSBub25jZQ==';            // the challenge (RFC sample)
const GUID = '258EAFA5-E914-47DA-95CA-C5AB0DC85B11'; // the fixed string
console.log(createHash('sha1').update(key + GUID).digest('base64'));
\end{code}

\begin{terminal}
$ node accept.js
s3pPLMBiTxaQ9kYGzzhZRbK+xOo=
\end{terminal}

The output is the \texttt{Accept} in the figure. That is all the code shows:
the answer is not magic, anyone who knows the fixed string can compute it.

The RFC has a second protection, and this one came from an attack. In 2010,
Lin-Shung Huang, Eric Chen, Adam Barth, Eric Rescorla and Collin Jackson showed
that, with the early drafts of the protocol, a malicious page could write
WebSocket messages whose bytes looked exactly like an ordinary HTTP request.
Some old proxies sitting between users and the internet didn't understand
WebSocket, read those bytes as a real request, and stored the attacker's content
in their cache under the address of a real site. Every other user behind that
proxy would then get the attacker's content when they visited that site. The
fix is called \emph{masking}: the browser scrambles every message it sends with
four random bytes, and sends the four bytes along so the server can unscramble
it. The page never sees those bytes, so it can no longer choose what goes
on the wire.

\begin{concept}{WebSocket}
A connection between a browser and a server that starts as an HTTP request and
then carries messages in both directions, each side sending whenever it wants,
until one side closes it.

\textbf{What it fixes:} the server can speak first, and each message costs a
few bytes of framing instead of a full request.

\textbf{What it costs:} the connection is state. The server has to keep it,
notice when it dies, and decide what to do when the other end reads slower
than it writes. HTTP gave you all of that for free by forgetting the client
after each answer.
\end{concept}

\section{The other way: Server-Sent Events}

Server-Sent Events would also carry this stream, and with less machinery. There
is no handshake and no change of protocol. The browser makes an ordinary HTTP
request, and the server answers with an ordinary HTTP response that never ends.
Each event is a few lines of text followed by a blank line. Here is a server
that sends two events, as \texttt{curl} sees it:

\begin{terminal}
$ curl -i http://localhost:8092/
HTTP/1.1 200 OK
Content-Type: text/event-stream
Date: Wed, 18 Nov 2026 16:38:10 GMT
Connection: keep-alive
Keep-Alive: timeout=5
Transfer-Encoding: chunked

id: 1
data: tick 1

id: 2
data: tick 2
\end{terminal}

The \texttt{data} line is the message. The \texttt{id} line is what makes the
format worth knowing. When the connection drops, the browser opens it again by
itself, without a line of code on the page, and the new request carries a
header, \texttt{Last-Event-ID}, with the last ID it received. The server reads
it and continues from there. In a test where the server dropped the connection
after every two events, the browser came back half a second later with
\texttt{Last-Event-ID: 2}, and the server went on from \texttt{tick 3}. The
half second was the server's choice: a \texttt{retry} line in the stream tells
the browser how long to wait. With a WebSocket, all of that (noticing the drop,
waiting, reconnecting, saying where you stopped) is code you write.

For data that only flows from the server to the browser, Server-Sent Events are
the simpler tool, and they are underused. The reason is history, not merit.
Internet Explorer never implemented them, and neither did the first version of
Edge. Microsoft's browser only got them in January 2020, when Edge was rebuilt
on Chromium. For ten years, any site that had to work in Internet Explorer used
a WebSocket library instead, and WebSocket became the default answer to ``how do
I push to the browser''.

The other weakness was a number: six connections. To see where it came from,
you have to know what changed between the versions of HTTP.

\begin{history}{1991--2022 \textperiodcentered{} HTTP, version by version}
The first HTTP, which Tim Berners-Lee wrote down in 1991 and which was later
called HTTP/0.9, had one command. The browser opened a connection, sent
\texttt{GET} and a path, received the page, and the server closed the
connection.

HTTP/1.0 (RFC~1945, May 1996) added headers, status codes and content types,
everything that lets an answer say what it is. It still used one connection per
request, and opening a TCP connection costs a round trip before the first byte
of the request goes out.

HTTP/1.1 (RFC~2068 in January 1997, revised as RFC~2616 in June 1999) kept the
connection open after the answer so the next request could reuse it. But on one
connection, requests still waited in line: the second answer could only start
after the first had finished. The standard allowed sending several requests
without waiting (\emph{pipelining}), but enough proxies and servers got it
wrong that browsers left it switched off. So browsers opened several connections
to the same site in parallel. RFC~2616 asked for no more than two. In 2008,
Internet Explorer~8 and Firefox~3 went to six, and the other browsers followed.
For a decade, web developers fought that number: they glued icons into one image
(\emph{sprites}), concatenated scripts into one file, and spread files over
several hostnames to get six more connections from each.

In November 2009, Mike Belshe and Roberto Peon, at Google, announced SPDY: many
requests on one connection at the same time, each cut into small frames
labeled with the request they belong to, so a slow answer no longer blocks the
others. Chrome shipped it in 2010. The IETF took SPDY as the starting point for
the next version, published as HTTP/2, RFC~7540, in May 2015. There was never
an HTTP/1.2.

HTTP/2 still runs on TCP, which delivers bytes strictly in order, so one lost
packet holds up every request on the connection until it is sent again.
HTTP/3 (RFC~9114, June 2022) keeps the same idea on top of QUIC, a transport
that Google also started, where a lost packet only holds up its own request.
\end{history}

\begin{concept}{Multiplexing}
Many conversations on one connection at the same time. Each one is cut into
small pieces, every piece is labeled with the conversation it belongs to, and
the other end puts each conversation back together.

\textbf{What it fixes:} the wait in line. A slow answer no longer blocks the
ones behind it, and there's no reason to open several connections to the same
place.

\textbf{What it costs:} everything rides on one connection. On TCP, a lost
packet stops all the conversations at once, which is what HTTP/3 fixed.
\end{concept}

\begin{etymology}
\emph{Multiplex} is Latin for ``many-fold''. Telegraph engineers used it in the
1870s for sending several messages over one wire at once; Thomas Edison's
\emph{quadruplex} of 1874 sent four. The problem was the same as HTTP's: wire
was expensive, so you made each one carry more.
\end{etymology}

Now the six connections make sense, and so does the trouble they gave
Server-Sent Events. In HTTP/1.1, an event stream is a response that never ends,
so it holds one of the six connections forever. The limit is per browser and per
site, shared by every tab. Open the operator screen in six tabs and all six
connections are taken. The seventh tab can't load anything from that server,
not even its own page, and it sits there loading until another tab closes.

\begin{figure}[htbp]
\centering
\begin{tikzpicture}[x=1cm, y=1cm]
  \node[hf/note, font=\sffamily\footnotesize\bfseries] at (2.5,4.2) {HTTP/1.1};
  \node[hf/reader, minimum width=14mm, minimum height=38mm] at (0,1.8) {browser};
  \node[hf/reader, minimum width=14mm, minimum height=38mm] at (5,1.8) {server};
  \foreach \y in {3.3,2.85,2.4,1.95,1.5,1.05} \draw[hf/push] (4.3,\y) -- (0.7,\y);
  \draw[hf/pain, dashed] (0.7,0.35) -- (2.2,0.35);
  \node[hf/note, anchor=west, text=hfAlert] at (2.25,0.35) {tab 7 waits};
  \node[hf/note, text width=50mm] at (2.5,-0.75) {six connections, each one held by a stream that never ends};
  \begin{scope}[xshift=7.5cm]
    \node[hf/note, font=\sffamily\footnotesize\bfseries] at (2.5,4.2) {HTTP/2};
    \node[hf/reader, minimum width=14mm, minimum height=38mm] at (0,1.8) {browser};
    \node[hf/reader, minimum width=14mm, minimum height=38mm] at (5,1.8) {server};
    \draw[black!40, line width=0.8pt, rounded corners=8pt] (0.7,0.55) rectangle (4.3,3.6);
    \foreach \y in {3.25,2.85,2.45,2.05,1.65,1.25,0.85} \draw[hf/push] (4.3,\y) -- (0.7,\y);
    \node[hf/note, text width=50mm] at (2.5,-0.75) {one connection, seven streams inside it};
  \end{scope}
\end{tikzpicture}
\caption{Seven tabs of the same screen, each with an event stream. In HTTP/1.1
every stream takes a whole connection, and the browser allows six per site. In
HTTP/2 the streams share one connection.}
\end{figure}

HTTP/2 removes the problem. Each event stream becomes one stream among many
inside a single connection, and the browser and the server agree on how many
can be open at once, usually a hundred. One catch survives: browsers speak
HTTP/2 only over HTTPS. On a plain \texttt{http://localhost}, which is where a
developer spends the day, it's HTTP/1.1 and the limit of six is back.

So the choice is simpler than its reputation. For a stream that only goes from
the server to the browser, use Server-Sent Events, served over HTTP/2. WebSocket
is the choice when the browser will also talk on the same connection, or when
the framework you use has better support for it.

Back to the screen. The chapter started with one question: how does the
backend get the picture of the plant to the browser without the browser asking
for it? Now there is an answer. The browser opens a connection once and keeps it
open, as a WebSocket or as an event stream, and the backend writes into it
whenever it has something new. The rest of this chapter uses WebSocket, but
almost everything in it holds for Server-Sent Events too.

An open connection isn't free, though. With plain HTTP, the server answered and
forgot the client. Now it has a list of connections it has to look after, and
the rest of the chapter is about what that asks of it: sending to many browsers
at once, a browser that stops reading, a browser that disappears without a word,
telling the screen that the plant has failed, agreeing on what a message looks
like, and testing all of it without a plant.

\section{One timer, many clients}

Start with the sending. With polling, every request from every tab made the
server read the plant and build an answer, so ten tabs meant ten times the work.
With open connections, the server has one timer instead. Every half second it
builds the picture once, turns it into text once, and writes that same text
into every open connection. Two tabs show the same moment of the plant, and the
work no longer depends on how many people are looking.

That loop has one weak point, and it is the subject of the rest of this section.
The server writes into each connection and moves on. It never checks whether the
browser on the other end is taking the messages out. Usually it is. But in a
control room some browser eventually isn't: a laptop on bad Wi-Fi, a browser
frozen by some other tab, a tablet left asleep on a desk. The connection is
still open, nobody closed it, and the server keeps writing into it twice a
second.

Those messages don't vanish. They wait, in the server's memory, and nothing
ever throws them away. One frozen tablet is enough to make the server's memory
grow, slowly, until somebody restarts it. The fix is a few lines: before
writing to a connection, look at how much is still waiting there, and if it's
too much, skip this picture for that browser. What follows shows why the
messages pile up, measures it, and explains why skipping is the right fix for a
screen.

\begin{question}{Doesn't the server just send and forget?}
That's how radio works: the station broadcasts, and whoever had the set on
heard it. A WebSocket doesn't work like that, because underneath it is a TCP
connection, and TCP makes one promise: every byte written arrives, complete and
in order, or the connection breaks. It never drops a byte quietly. To keep
that promise, the sending side has to hold on to every byte until the other
side confirms it has it. A browser that isn't reading confirms nothing, so the
bytes stay.

It isn't one broadcast, either. Each browser has its own connection, and the
loop that ``sends to everyone'' is a \texttt{for} that writes a copy into each
connection, one at a time.

Send-and-forget does exist on the internet. In 1978, TCP was split in two: IP
to carry packets, and TCP on top of it for the promise. Danny Cohen, who was
sending voice over the network, had argued for the split, because in a
conversation a word that arrives late is worse than a word that's lost. David
P.~Reed designed UDP for programs like his, packets with no promise, and it
was published as RFC~768 in August 1980. Video calls and games still run on it.
A page that pushes data to a screen, like this one, runs on TCP.
\end{question}

When the server calls \texttt{send()}, the call returns at once, but the
message hasn't arrived anywhere. To see where it is, you need the pieces a
message passes through on its way from the backend to the screen, and who
writes each one. Figure~\ref{fig:buffers} shows them.

\begin{figure}[htbp]
\centering
\begin{tikzpicture}[x=1cm, y=1cm]
  \tikzset{col/.style={minimum width=44mm, minimum height=10mm}}
  % server machine
  \node[hf/note, font=\sffamily\footnotesize\bfseries] at (0,1.75) {server machine};
  \node[hf/reader, col] (code) at (0,0) {your code\\{\footnotesize the timer loop: build, check, \texttt{send()}}};
  \node[hf/bus, col] (lib) at (0,-1.6) {WebSocket library (\texttt{ws})\\{\footnotesize makes frames, reports \texttt{bufferedAmount}}};
  \node[hf/bus, col, draw=hfAlert, line width=1.4pt] (sock) at (0,-3.2) {Node's socket\\{\footnotesize a queue with no limit}};
  \node[hf/source, col] (kout) at (0,-5.4) {operating system: TCP\\{\footnotesize send buffer, fixed size}};
  \begin{scope}[on background layer]
    \node[draw=black!35, dashed, rounded corners=3pt, inner sep=3mm,
      fit=(code)(lib)(sock), label={[hf/note, anchor=south]north:Node process, server memory}] {};
  \end{scope}
  \draw[hf/push] (code) -- (lib);
  \draw[hf/push] (lib) -- (sock);
  \draw[hf/push] (sock) -- (kout);
  \draw[hf/ask, dashed] (sock.west) to[bend left=50]
    node[hf/note, left, text=hfReader, align=right] {reads\\\texttt{bufferedAmount},\\skips if too big} (code.west);
  \node[hf/note, text=hfAlert, anchor=west, align=left] at ([xshift=4mm]sock.east) {grows with\\every tick};
  % client machine
  \node[hf/note, font=\sffamily\footnotesize\bfseries] at (6.9,1.75) {client machine};
  \node[hf/reader, col] (page) at (6.9,0) {the page\\{\footnotesize draws the screen}};
  \node[hf/reader, col] (br) at (6.9,-1.6) {browser\\{\footnotesize WebSocket, \texttt{onmessage}}};
  \node[hf/source, col] (kin) at (6.9,-5.4) {operating system: TCP\\{\footnotesize receive buffer, fixed size}};
  \draw[hf/push] (br) -- (page);
  \draw[hf/push] (kin) -- (br);
  \draw[hf/pain] (6.65,-3.75) -- (7.15,-3.25) (6.65,-3.25) -- (7.15,-3.75);
  \node[hf/note, text=hfAlert, anchor=west] at (7.3,-3.5) {stops reading};
  % the wire
  \draw[hf/push] ([yshift=2mm]kout.east) -- node[hf/note, above] {network} ([yshift=2mm]kin.west);
  \draw[hf/pain, dashed, -{Stealth[length=2.2mm]}] ([yshift=-2mm]kin.west) --
    node[hf/note, below, text=hfAlert, align=center] {full: stop sending\\(the TCP window)} ([yshift=-2mm]kout.east);
\end{tikzpicture}
\caption{From the backend to the screen. Three pieces live inside the Node
process; TCP and its two buffers belong to the operating system on each
machine. When the browser stops reading, the two TCP buffers fill and stop, and
the messages pile up in Node's socket, in red. The check that skips a picture
is in your code, before \texttt{send()}.}
\label{fig:buffers}
\end{figure}

On the server, three pieces live inside the Node process, in its memory:

\begin{itemize}
\item \textbf{Your code}, the timer loop. It builds the picture and calls
  \texttt{send()} once for each connection. The web framework, Express or
  Fastify, isn't in this path at all: it answered the HTTP request that opened
  the connection and handed the connection over.
\item \textbf{The WebSocket library}, usually \texttt{ws}. It wraps the
  message in a WebSocket frame and passes it down. It also tells you how big
  the queue below it is, as \texttt{bufferedAmount}.
\item \textbf{Node's socket}, one per connection. This is where bytes wait
  until the operating system accepts them, and this queue has no maximum.
\end{itemize}

Below Node is the operating system. TCP isn't written in Node, or in the
library: it's part of the operating system, on both machines, and Node only
asks it to send some bytes. The operating system keeps a \emph{send buffer} of
fixed size for each connection, sends from it, and only deletes a byte once the
other side confirms it arrived. On the client, its operating system keeps a
\emph{receive buffer}, also fixed, where bytes wait until the browser reads
them. The browser then unwraps the frame and calls the page's
\texttt{onmessage}, and the page draws.

Now the browser stops reading. The chain fills from the end. The receive
buffer on the client fills, and the client's operating system tells the
server's: no room, stop sending. The server's send buffer fills next, because
nothing leaves it. From then on the operating system accepts nothing more from
Node, and every new message stays in Node's socket. The two TCP buffers have a
fixed size and stop growing there. Node's socket doesn't. It's the red box in
the figure, and it grows with every tick, in the server's memory, for as long as
that browser stays connected.

The skip goes in your code, the only piece in the figure you write. Before
calling \texttt{send()} for a connection, the loop looks at
\texttt{bufferedAmount}. If it's above a limit, the loop doesn't call
\texttt{send()} for that connection on this tick. Nothing new enters the red
box, and the browser simply gets the latest picture once it starts reading
again.

The ``no room, stop sending'' between the two machines isn't something Node
does. It's built into TCP, and it's fifty years old.

\begin{history}{1974--2015 \textperiodcentered{} Flow control}
In 1974, Vint Cerf and Bob Kahn described TCP with a \emph{window}: the receiver
tells the sender how many bytes it still has room for, and the sender never
sends more. A slow reader slows the writer down, all the way back through the
network. In software the same idea came back under the name
\textbf{backpressure}. Node's streams had it from the start: \texttt{write()}
returns \texttt{false} when the stream's buffer is full, and a \texttt{drain}
event says when to continue. From 2013 to 2015, engineers from Netflix,
Typesafe (now Lightbend), Pivotal and others wrote the Reactive Streams specification so that
asynchronous libraries on the JVM could pass the signal to one another, and
Java~9 adopted it in 2017.
\end{history}

TCP's ``stop'' reaches the server, but the WebSocket's \texttt{send()} in Node
doesn't pass it on to your code. It never refuses a message and never says
``wait''. It queues and returns. The
only trace of a reader that has stopped is a number on the connection,
\texttt{bufferedAmount}: how many bytes are sitting in Node's socket, the red
box in Figure~\ref{fig:buffers}. A server that never looks at it will keep
everything it ever tried to send to a stalled client.

To see it happen, here is a server that sends 16~KB every millisecond to every
client and prints the size of each queue once a second. It also takes an
optional limit: a client whose queue is already longer than that is skipped on
this tick.

\begin{code}
import { WebSocketServer } from 'ws';

const LIMIT = Number(process.env.LIMIT ?? Infinity); // bytes
const wss = new WebSocketServer({ port: 8090 });
const message = 'x'.repeat(16 * 1024);

setInterval(() => {
  for (const client of wss.clients) {
    if (client.bufferedAmount > LIMIT) continue; // skip this tick
    client.send(message);
  }
}, 1);

setInterval(() => {
  for (const client of wss.clients) console.log('buffered:', client.bufferedAmount);
}, 1000);
\end{code}

On the other end is the stalled client: a small program that opens a
connection, sends the WebSocket handshake shown earlier, and then never
reads a single byte. The server ran twice, on the left without a limit and on
the right with a limit of 64~KiB. Each line is one second, and the number is
the size of that client's queue in bytes:

\begin{terminal}
$ node flood.js                    $ LIMIT=65536 node flood.js
buffered: 0                        buffered: 0
buffered: 0                        buffered: 0
buffered: 491640                   buffered: 65552
buffered: 1622412                  buffered: 65552
buffered: 2769572                  buffered: 65552
buffered: 4211716                  buffered: 65552
buffered: 5948844                  buffered: 65552
buffered: 7096004                  buffered: 65552
\end{terminal}

\begin{figure}[htbp]
\centering
\begin{tikzpicture}[x=1.2cm, y=0.6cm]
  \draw[black!40, -{Stealth[length=2mm]}] (0,0) -- (8.6,0) node[hf/note, right] {seconds};
  \draw[black!40, -{Stealth[length=2mm]}] (0,0) -- (0,7.8) node[hf/note, above] {MB queued};
  \foreach \t in {1,...,8} \draw[black!40] (\t,0.08) -- (\t,-0.08) node[hf/note, below] {\t};
  \foreach \m in {2,4,6} \draw[black!40] (0.08,\m) -- (-0.08,\m) node[hf/note, left] {\m};
  \draw[hfAlert, line width=1pt] plot coordinates {(1,0) (2,0) (3,0.49) (4,1.62) (5,2.77) (6,4.21) (7,5.95) (8,7.10)};
  \foreach \t/\v in {1/0, 2/0, 3/0.49, 4/1.62, 5/2.77, 6/4.21, 7/5.95, 8/7.10} \fill[hfAlert] (\t,\v) circle (1.8pt);
  \node[hf/note, text=hfAlert, anchor=west] at (6.2,6.6) {no limit};
  \draw[hfReader, line width=1pt] plot coordinates {(1,0) (2,0) (3,0.066) (4,0.066) (5,0.066) (6,0.066) (7,0.066) (8,0.066)};
  \node[hf/note, text=hfReader, anchor=south west] at (5.2,0.2) {skip above 64 KiB};
\end{tikzpicture}
\caption{The two runs above. Without a limit, a client that doesn't read costs
the server more than a megabyte every second. With the limit, the queue stops at
64 KiB.}
\end{figure}

The two seconds at zero don't mean there was no problem. They are the operating
system's buffers, on both ends, filling up before Node's own queue starts to
grow. And the experiment is exaggerated on purpose: at a real rate of one
picture every half second, the same stalled client would take hours to cost
anything worth noticing. The leak is real and slow, and the fast loop only
compresses time.

\begin{concept}{Backpressure}
The signal a slow consumer sends back to a fast producer: stop, or slow down.
When the producer can't slow down, it has to choose what to do with what it
can't deliver: block, keep it in memory, or throw it away.

\textbf{What it fixes:} one slow reader can't grow the producer's memory until
the process dies.

\textbf{What it costs:} a decision about which data is allowed to be lost, and
it has to be made for each stream separately.
\end{concept}

For a live picture of the plant the decision is easy. A screen only needs the
latest picture, and an older one that arrives late is worse than none. So a
client whose queue is too long skips ticks until it drains, and then gets the
current picture. This is the same reasoning as the queue of one for the screen
in Chapter~2. A stream that is being recorded is the opposite case: skipping is
data loss, and the reader has to be able to tell that it missed something,
which takes a sequence number in every message.

\section{Is anyone there?}

Skipping protects the server's memory from a slow client. A dead client is a
different problem. If a laptop is closed or a cable is pulled, no packet
announces it. The server's socket stays open and looks healthy, and the server
keeps queueing data for it.

\begin{history}{1989 \textperiodcentered{} Two hours of silence}
TCP has a keep-alive of its own, but RFC~1122, the 1989 rules for Internet
hosts, made it optional and said that if it's used, the default interval must be
at least two hours. The worry at the time was cost: a probe packet on a link
billed by the packet, to check on a connection nobody was using. Most operating
systems still default to two hours. That's why protocols that need to know
sooner, from OPC UA's keep-alive in Chapter~2 to WebSocket, carry their own.
\end{history}

WebSocket's own check is made of frames, so first, what a frame is. After the
handshake, everything that travels on the connection is a \emph{frame}: a header
of two bytes or more, then the content. The first byte says what kind of frame
it is, in its last four bits, a number called the \emph{opcode}: 1 is text, 2 is
binary, 8 is close, 9 is \emph{ping} and 10 is \emph{pong}. The second byte
gives the length, and says whether the frame is masked. Here are three frames
from a real connection, in hexadecimal, as a client that reads raw bytes saw
them:

\begin{terminal}
client got: 81 02 68 69
client got: 89 00
client sends: 8a 80 12 34 56 78
\end{terminal}

The first is a text frame. In \texttt{81}, the 8 says this is the last piece of
its message and the 1 says text; \texttt{02} is the length, and \texttt{68 69}
is ``hi''. The second, \texttt{89 00}, is a ping with nothing inside: two bytes
in all. The third is the answer going back. In \texttt{8a}, the \texttt{a} is
10, a pong; \texttt{80} says masked, length zero; and the last four bytes are
the mask, which every frame from a client carries, as the masking story earlier
in the chapter explained. In this test the client answered by hand. A browser
does it by itself.

Ping and pong are \emph{control frames}. They belong to the connection, not to
the application. In Figure~\ref{fig:buffers}, a ping is made by the WebSocket
library when your code calls \texttt{ping()}, and travels down the same path as
a picture. On the client it stops one floor early: the browser answers it with a
pong on its own, and the page never finds out. The browser's WebSocket API
doesn't even have an event for pings. Back on the server, the library tells your
code with a \texttt{pong} event.

That's enough to find a dead client. The server checks its clients in rounds,
on a slow timer, every thirty seconds or so. In each round it looks at every
connection. One that has answered since the last round is marked as unknown and
pinged again. One that is still unknown never answered the last ping, and gets
closed.

\begin{figure}[htbp]
\centering
\begin{tikzpicture}[x=1cm, y=-0.75cm]
  \node[hf/reader, minimum width=17mm] at (0,0) {page A};
  \node[hf/reader, minimum width=17mm] at (2.8,0) {browser A};
  \node[hf/reader, minimum width=17mm] at (6.1,0) {server};
  \node[hf/reader, minimum width=17mm] at (9.4,0) {laptop B};
  \draw[hf/lifeline] (0,0.7) -- (0,7.4);
  \draw[hf/lifeline] (2.8,0.7) -- (2.8,7.4);
  \draw[hf/lifeline] (6.1,0.7) -- (6.1,7.4);
  \draw[hf/lifeline] (9.4,0.7) -- (9.4,6.3);
  % an ordinary message goes all the way to the page
  \draw[hf/push] (6.1,1.3) -- node[hf/note, above] {picture (text, 1)} (2.8,1.3);
  \draw[hf/push] (2.8,1.7) -- node[hf/note, above] {\texttt{onmessage}} (0,1.7);
  % round 1
  \draw[black!35, dashed] (-0.9,2.4) -- (10.3,2.4);
  \node[hf/note, anchor=east] at (-1.0,2.4) {round 1};
  \draw[hf/push] (6.1,3.2) -- node[hf/note, above] {ping (9)} (2.8,3.2);
  \draw[hf/push] (6.1,3.6) -- node[hf/note, above] {ping (9)} (9.4,3.6);
  \draw[hf/answer] (2.8,4.1) -- node[hf/note, above] {pong (10)} (6.1,4.1);
  \node[hf/note, text=hfAlert, anchor=east] at (9.3,4.4) {no answer};
  \node[hf/note, fill=white, align=center] at (0,3.5) {never sees\\the ping};
  % round 2
  \draw[black!35, dashed] (-0.9,5.0) -- (10.3,5.0);
  \node[hf/note, anchor=east, align=right] at (-1.0,5.0) {round 2,\\30\,s later};
  \draw[hf/push] (6.1,5.8) -- node[hf/note, above] {ping (9)} (2.8,5.8);
  \draw[hf/pain, -{Stealth[length=2.2mm]}] (6.1,6.3) -- node[hf/note, above, text=hfAlert] {still unknown: close} (9.4,6.3);
  \draw[hf/answer] (2.8,6.8) -- node[hf/note, above] {pong (10)} (6.1,6.8);
\end{tikzpicture}
\caption{Two rounds of the check. A picture goes all the way to the page; a ping
stops at the browser, which answers it by itself. Browser A answers every round
and stays. Laptop B, closed, doesn't answer round 1, and in round 2 the server
closes its connection.}
\label{fig:heartbeat}
\end{figure}

In Node, with \texttt{ws}, the whole check is a timer and a flag on each
connection:

\begin{code}
setInterval(() => {
  for (const ws of wss.clients) {
    if (!ws.isAlive) { ws.terminate(); continue; } // didn't answer last round
    ws.isAlive = false;                             // unknown until it answers
    ws.ping();
  }
}, 30_000);

wss.on('connection', (ws) => {
  ws.isAlive = true;
  ws.on('pong', () => { ws.isAlive = true; });
});
\end{code}

In a test with rounds of one second instead of thirty, there were two clients:
A, a real WebSocket client, and B, a program that did the handshake and then
went silent. A answered at 1.0, 2.0 and 3.0 seconds. B was closed at 2.0 seconds, in
the second round after it connected, which is the worst case: a dead client
lasts at most two rounds.

\section{A failure is data}

The last two sections were about the browser's side of the backend: a browser
that stops reading, a browser that disappears. The backend has another side,
the plant, and that side fails too. Every half second, the timer reads the plant
over OPC UA, as in Chapter~2, and pushes the picture. Sooner or later one of
those reads will fail: the plant's server restarts, the network drops for a few
seconds, the plant's server crashes. This section is about what the backend,
and the operator's screen, should do then.

The first instinct is the one most code has for something going wrong: the read
throws an error, \texttt{connection closed}. In a function called by a timer,
that's a bad idea, because nobody is there to catch it. Your code didn't call
the function, Node's timer did, and there's no \texttt{try} around it. The
program below is that situation reduced to three lines: a timer whose function
fails on its first tick, the way a read of a plant that's gone would fail. It
prints \texttt{streaming} when it starts:

\begin{code}
setInterval(async () => {
  throw new Error('connection closed');
}, 500);
console.log('streaming');
\end{code}

\begin{terminal}
$ node crash.js
streaming
Error: connection closed
    at Timeout._onTimeout (file:///.../crash.js:2:9)
$ echo $?
1
\end{terminal}

Half a second after it started, the process was dead. \texttt{echo \$?} prints
the exit code of the last program, and 1 means it ended in failure. In the real
backend, that's one hiccup in the plant's network killing the whole server, and
with it the connection of every operator's screen. The function is
\texttt{async}, so the error doesn't come out as a plain exception but as a
rejected promise.

\begin{concept}{Promise}
An object that stands for a result that isn't ready yet. A function that has to
wait (for the network, a disk, a timer) returns a promise at once, and later the
promise \emph{resolves} with the value or \emph{rejects} with an error. Whoever
holds it can wait for it with \texttt{await}, or attach handlers with
\texttt{.then()} and \texttt{.catch()}. A function marked \texttt{async} always
returns a promise: what it returns becomes the value, and what it throws becomes
the rejection. When a promise settles, the code waiting for it doesn't run on
the spot: it goes into the microtask queue from Chapter~3, and runs as soon as
the current callback ends.

\textbf{What it fixes:} code that waits reads in order, top to bottom, and a
failure travels back to whoever is waiting, the way an exception does in code
that doesn't wait.

\textbf{What it costs:} the error only reaches someone who is waiting. If nobody
awaits the promise or attaches a \texttt{.catch()}, the rejection has nowhere to
go. That's the timer above: \texttt{setInterval} calls the function and ignores
the promise it returns.
\end{concept}

Nobody was waiting for this one, and for years Node handled that case
differently.

\begin{history}{1976--2020 \textperiodcentered{} Promises, and the ones nobody waited for}
The idea is older than JavaScript. Daniel Friedman and David Wise used the word
\emph{promise} in 1976, and Henry Baker and Carl Hewitt called the same thing a
\emph{future} in 1977: a placeholder for a value that is still being computed
somewhere else. In 1988, Barbara Liskov and Liuba Shrira used promises for calls
between machines.

JavaScript got there through pain. Node's API was built on callbacks: you passed
a function to be called with the error or the result when the work was done.
Every step that waited nested the code one level deeper, ``callback hell'', and
every level had to pass errors along by hand. Reading two files, writing them
together into a third and reading that back takes four steps, and four steps
make a pyramid:

\begin{center}
\setbox0\hbox{\ttfamily\scriptsize 0}\edef\hfcw{\the\wd0}%
\begin{tikzpicture}[x=1cm, y=1cm]
  \def\q{\textquotesingle}
  \tikzset{ln/.style={anchor=west, inner sep=0pt, font=\ttfamily\scriptsize, text=black!80},
    er/.style={ln, text=hfAlert}}
  \def\lh{0.37}
  % the pyramid: the empty space the indentation leaves
  \fill[hfAlert!8] (0,0.15) -- (8*\hfcw,{-8*\lh}) -- (0,{-12*\lh-0.15}) -- cycle;
  \node[ln] at (0,0) {fs.readFile(\q a.txt\q, \q utf8\q, (err, a) => \{};
  \node[er] at (2*\hfcw,{-1*\lh}) {if (err) return console.error(err);};
  \node[ln] at (2*\hfcw,{-2*\lh}) {fs.readFile(\q b.txt\q, \q utf8\q, (err, b) => \{};
  \node[er] at (4*\hfcw,{-3*\lh}) {if (err) return console.error(err);};
  \node[ln] at (4*\hfcw,{-4*\lh}) {fs.writeFile(\q ab.txt\q, a + b, (err) => \{};
  \node[er] at (6*\hfcw,{-5*\lh}) {if (err) return console.error(err);};
  \node[ln] at (6*\hfcw,{-6*\lh}) {fs.readFile(\q ab.txt\q, \q utf8\q, (err, ab) => \{};
  \node[er] at (8*\hfcw,{-7*\lh}) {if (err) return console.error(err);};
  \node[ln] at (8*\hfcw,{-8*\lh}) {console.log(ab);};
  \node[ln] at (6*\hfcw,{-9*\lh}) {\});};
  \node[ln] at (4*\hfcw,{-10*\lh}) {\});};
  \node[ln] at (2*\hfcw,{-11*\lh}) {\});};
  \node[ln] at (0,{-12*\lh}) {\});};
  \node[hf/note, text=hfAlert, anchor=north west, align=left] at (0,{-12*\lh-0.3})
    {callbacks: one level deeper per step,\\and the error handled by hand on each};
  % the same four steps, flat
  \begin{scope}[xshift=7.5cm, yshift={-9*\lh cm}]
    \node[ln] at (0,0) {const a = await readFile(\q a.txt\q, \q utf8\q);};
    \node[ln] at (0,{-1*\lh}) {const b = await readFile(\q b.txt\q, \q utf8\q);};
    \node[ln] at (0,{-2*\lh}) {await writeFile(\q ab.txt\q, a + b);};
    \node[ln] at (0,{-3*\lh}) {console.log(await readFile(\q ab.txt\q, \q utf8\q));};
    \node[hf/note, text=hfReader, anchor=north west, align=left] at (0,{-3*\lh-0.3})
      {\texttt{await}: one line per step, and a failure\\travels by itself to whoever is waiting};
  \end{scope}
\end{tikzpicture}
\end{center}

The two programs do the same work: with \texttt{one} in \texttt{a.txt} and
\texttt{two} in \texttt{b.txt}, both print \texttt{one two}. Only the shape
changes. Libraries filled the gap: Kris
Zyp's Promises/A for CommonJS in 2009, jQuery's Deferred in 2011, and
Promises/A+ in 2012, a short specification that made the libraries work with
each other. ES2015 put \texttt{Promise} in the language, and ES2017 added
\texttt{async} and \texttt{await}.

When a promise fails and nobody is waiting for it, the failure is an
\emph{unhandled rejection}. For years, Node only printed a warning about it and
went on running, in a state nobody had planned for. From Node~7, in 2016, the
warning said that a future version would end the process instead. Node~15, in October 2020, made
that the default: an unhandled rejection now kills the process, the same as an
uncaught exception always had.
\end{history}

Catching the error and writing it to a log would keep the process alive, but
it fixes only half the problem. The browser would still know nothing: the
pictures would stop coming, and the screen would go on showing the last numbers
it got, as if the plant were steady. That's where the title of this section
comes from. A lost plant isn't an accident; it's a state the system will be in,
sooner or later, so it belongs in the data, not in an error. Each message to
the browser says whether the backend is connected to the plant or
reconnecting, and while it's reconnecting, the picture is empty. The backend retries with the exponential
backoff from Chapter~2, forever, and the browser sees the state change as it
happens.

This is the rule about data quality from Chapter~2, applied one hop further:
the browser can't tell a steady plant from a dead connection unless somebody
tells it. A screen that keeps showing the last numbers after the plant is gone
is lying.

\section{One definition of the message}

The last section changed the message. It used to be just the 53 numbers of the
plant; now it also says whether the backend is connected to the plant, and the
numbers can be missing. A change like that has to happen in two programs at
once: the backend, which writes the message, and the browser, which reads it.
They are different programs, possibly in different repositories. If each one
writes down the shape of the message on its own, sooner or later the two will
disagree, and nothing will warn anyone. This section is about writing that shape
down once.

There is also a quieter problem, at the point where values from the plant
enter the backend:

\begin{code}
const raw = ['2696.1', 'n/a', '74.2'];
console.log(JSON.stringify({ values: raw.map(Number) }));
\end{code}

\begin{terminal}
$ node nan.js
{"values":[2696.1,null,74.2]}
\end{terminal}

\texttt{Number('n/a')} is \texttt{NaN}. JSON has no way to write \texttt{NaN},
so \texttt{JSON.stringify} quietly writes \texttt{null}, and the screen receives
a hole in the middle of a list of numbers with no error anywhere. The type
system can't help: \texttt{NaN} is a number.

\begin{history}{1987--2009 \textperiodcentered{} Interface definitions}
Programs that talk to each other have been describing their messages in a
separate, neutral language for a long time. Sun's RPC used XDR, from 1987, with
a tool that generated the C code for both ends. CORBA, in 1991, had its IDL.
Google released Protocol Buffers in 2008. When JSON took over the web, it came
without any of that: Douglas Crockford's format, from 2001, describes values,
not their shapes. In 2009, Kris Zyp published the first draft of
\textbf{JSON Schema}, a way to describe the shape of JSON in JSON. It never
became an RFC, its current version is called 2020-12, and it's what most tools
that need a language-neutral contract for JSON read and write.
\end{history}

With a library like zod (Chapter~3), the shape is written once, in code, and
three things come out of it: the TypeScript type, a check that runs on real
data, and a JSON Schema for programs written elsewhere.

\begin{code}
import { z } from 'zod';

const Reading = z.object({
  status: z.enum(['connected', 'reconnecting']),
  values: z.array(z.number()).length(3).nullable(),
});

export type Reading = z.infer<typeof Reading>;
console.log(JSON.stringify(z.toJSONSchema(Reading), null, 2));
\end{code}

\begin{terminal}
$ node toschema.ts
{
  "$schema": "https://json-schema.org/draft/2020-12/schema",
  "type": "object",
  "properties": {
    "status": {
      "type": "string",
      "enum": [
        "connected",
        "reconnecting"
      ]
    },
    "values": {
      "anyOf": [
        {
          "minItems": 3,
          "maxItems": 3,
          "type": "array",
          "items": {
            "type": "number"
          }
        },
        {
          "type": "null"
        }
      ]
    }
  },
  "required": [
    "status",
    "values"
  ],
  "additionalProperties": false
}
\end{terminal}

\begin{figure}[htbp]
\centering
\begin{tikzpicture}
  \node[hf/bus, minimum width=34mm, minimum height=12mm] (s) {one schema};
  \node[hf/reader, right=20mm of s, yshift=14mm, minimum width=36mm] (t) {TypeScript type\\{\footnotesize checked while compiling}};
  \node[hf/reader, right=20mm of s, minimum width=36mm] (p) {runtime check\\{\footnotesize every value that comes in}};
  \node[hf/source, right=20mm of s, yshift=-14mm, minimum width=36mm] (j) {JSON Schema file\\{\footnotesize committed}};
  \node[hf/reader, right=12mm of j, minimum width=26mm] (w) {browser code\\{\footnotesize its own types}};
  \draw[hf/flow, black!60] (s.east) -- (t.west);
  \draw[hf/flow, black!60] (s.east) -- (p.west);
  \draw[hf/flow, black!60] (s.east) -- (j.west);
  \draw[hf/flow, black!60] (j) -- (w);
\end{tikzpicture}
\caption{The shape of the message is written once. The other program never
reads the backend's code, only the generated file.}
\end{figure}

\begin{concept}{Schema as contract}
The shape of a message between two programs is defined in one place, and every
other representation (types, runtime checks, documentation, code in another
language) is generated from it.

\textbf{What it fixes:} the two sides can't drift apart silently, and a bad
value is rejected where it enters instead of reaching the screen as a hole.

\textbf{What it costs:} a generation step, and a generated file that can fall
out of date. The usual guard is a test that regenerates it and fails if the
result differs from what's committed.
\end{concept}

There is one trap in validating where data enters. When a read from the plant
fails the check, it's tempting to send it down the same path as a lost
connection, and then the screen says ``reconnecting'' about a plant that is
perfectly reachable and sending something wrong. ``The source is gone'' and
``the source is sending garbage'' are different states, and an operator needs to
know which one is happening.

\section{A source you can swap}

The screen has to be built and tested by people who don't always have a plant
running. So the backend shouldn't decide by itself where the picture comes
from.

\begin{history}{2004--2007 \textperiodcentered{} Injection and doubles}
In January 2004, Martin Fowler wrote an article about the containers that were
then appearing in Java, and gave a name to what they did: \textbf{dependency
injection}. An object doesn't build the things it depends on; it receives them
from outside, so whoever builds it chooses. In 2007, Gerard Meszaros's
\emph{xUnit Test Patterns} sorted out the vocabulary for what tests put in
those places (dummies, stubs, fakes, mocks) under one name, the \textbf{test
double}, after the stunt double who stands in for an actor.
\end{history}

In a small backend, injection needs no container. The function that builds the
application receives the source of the picture as a parameter: a function that
returns the next message. The program decides, when it starts, which one to
pass. In production, it's the plant. In tests, it's a fake that returns invented
numbers instantly. On a developer's machine without a plant, it's a recording:
a file of past pictures played back in a loop.

A recording has one rule that matters. If it picks the row to serve by counting
calls, every extra reader of the source (a second tab, a test, a health check)
makes it skip rows. If it picks the row by the clock, any number of readers see
the same moment.

And a replay must say that it's a replay. A recording played back with the
status ``connected'' looks exactly like a live plant, and sooner or later
somebody leaves it switched on where it shouldn't be. Every message should say
where it came from.

\section{What it costs}

The backend now holds one session with the plant and an open connection with
every browser, and it is responsible for both. It has to notice when either end
dies, decide what a slow reader is allowed to lose, and say in every message
whether the data is live. A page that asked every half second needed none of
that. Each of those responsibilities is the price of the server being able to
speak first.

\section*{Further reading}

\begin{itemize}
  \item I.\ Fette and A.\ Melnikov, \emph{The WebSocket Protocol}, RFC~6455,
        2011, \url{https://www.rfc-editor.org/rfc/rfc6455}.
  \item WHATWG, \emph{HTML Living Standard}, section ``Server-sent events'',
        \url{https://html.spec.whatwg.org/multipage/server-sent-events.html}.
  \item J.~J.\ Garrett, ``Ajax: A New Approach to Web Applications'', 2005.
  \item L.-S.\ Huang, E.\ Chen, A.\ Barth, E.\ Rescorla and C.\ Jackson,
        ``Talking to Yourself for Fun and Profit'', 2010.
  \item V.\ Cerf and R.\ Kahn, ``A Protocol for Packet Network
        Intercommunication'', \emph{IEEE Transactions on Communications} 22(5),
        1974.
  \item R.\ Braden (ed.), \emph{Requirements for Internet Hosts --
        Communication Layers}, RFC~1122, 1989.
  \item Reactive Streams, \url{https://www.reactive-streams.org}.
  \item JSON Schema, \url{https://json-schema.org}.
  \item M.\ Fowler, ``Inversion of Control Containers and the Dependency
        Injection pattern'', 2004,
        \url{https://martinfowler.com/articles/injection.html}.
  \item G.\ Meszaros, \emph{xUnit Test Patterns}, Addison-Wesley, 2007.
\end{itemize}
